%%%%%%%%%%%%%%%%%%%%%%%%%%%%%%%%%%%%%%%%%%%%%%%%%%%%%%%%%%%%%%%%%%%%%%%%%%%%%%%%%%%%%%%%%%%%%%%%%%%%%%
%
%   Filename    : chapter_2.tex 
%
%   Description : This file will contain your review of related works.
%                 
%%%%%%%%%%%%%%%%%%%%%%%%%%%%%%%%%%%%%%%%%%%%%%%%%%%%%%%%%%%%%%%%%%%%%%%%%%%%%%%%%%%%%%%%%%%%%%%


%%%%%%%

\chapter{Related Works}
\label{sec:relatedworks}
\section{Cultural differences in emoji use}

Eastern and western cultures differed dramatically in many aspects, such as their languages, customs, cuisines, beliefs, and many more. Although many were all connected through the Internet, the way easterners and westerners communicated online and used emojis still had many differences. A study by Guntuku et al. (\citeyear{Guntuku2019}) found that easterners and westerners used different emojis to signify hunger. More western cultures tended to use pasta or meat, while eastern cultures tended to use noodles and rice bowls to signify food and hunger. The study also mentioned that Eastern cultures tended to use more emojis with faces than western users and that western users tended to employ a wider array of emojis than eastern users, who frequently used only a small amount of emojis. Another study by Togans et al. (\citeyear{TOGANS2021277}) shared some findings with the study mentioned above. One of them was that easterners used emojis in their texts more frequently than westerners and that the interdependent and interpersonal nature of easterners affected the way they used emojis compared to the independent nature of westerners. This study also found that eastern cultures were more likely to use emojis that fit the context and mood of the current conversation to avoid misunderstanding.

Meanwhile, the Philippines, similar to some other nations, had both eastern and western influences when it came to culture. A paper by Dolot \& Opina (\citeyear{Dolot2021}) went into detail about how Filipinos used graphicons such as images, videos, stickers, and emojis. The most frequently used graphicon was emojis. The study found that Filipinos had some behaviors similar to those mentioned in previous studies but also exhibited unique traits not discussed earlier, such as the way Filipinos tended to use multiple different emojis together, used emojis to grab people’s attention, and sometimes used emojis non-sensically to convey humor.

\section{\hl{Cultural Interpretation of Emojis}}

Different cultures interpret emojis differently, the cross-cultural communication study by Togans et. al. observed younger Japanese users, especially women, incorporate emoji heavily into messaging as part of “kawaii” culture (Dharma et. al., 2022). These findings suggest that the Philippines, which ranks among the highest countries using emojis outside the US, likely exhibits its own distinctive emoji norms (Kejriwal et. al., 2021). In fact, one study noted that the Philippines alongside Brazil and India as the top users of emoji outside the US. Table 1 summarizes the reported emoji usage statistics in Western versus Filipino and East Asian contexts.

%Table for Emoji Usage statistics in culture%
\begin{center}
	\begin{tabular}{m{7em} m{20em} m{5em}}
		Region/Culture & Usage/Statistic & Source \\ [0.5ex]
		\hline
		United States & 74\% of smartphone users reported using emojis in messaging apps (2013) & Nishimura(\citeyear{Nishimura2015}) \\
		\hline
		Japan & Emoji use is widespread, especially among younger women influenced by 'kawaii' culture & Nishimura(\citeyear{Nishimura2015})\\
		\hline
		East Asia & East Asian users employ emojis more frequently and judiciously than Western users & Togans et al (\citeyear{TOGANS2021277}) \\
		\hline
		Philippines & Among highest global emoji usage; identified as a top country for emoji use outside the US & Mayank Kejriwal (\citeyear{Kejriwal2021}) \\
	\end{tabular}
\end{center}

An example of cultural differences in emoji perception can be seen in the 2021 ESR valence ranking. The “\emoji{smiling-face}” (smiling-face) emoji is considered highly positive in Western cultures. However, in the Philippines, it can convey feelings of shyness or romantic excitement, known as “kilig,” which is unique to Filipino culture. Similarly, while the “\emoji{face-with-tears-of-joy}” (face with tears of joy) emoji is rated very positively by English-based sentiment models, Filipino Gen Z users often use the “\emoji{skull}” (skull) emoji to express sarcasm and mild humor. This demonstrates that cultural background and age significantly influence how emojis are interpreted. As a result, general sentiment models may misinterpret many expressions within the Filipino language.

Western users often interpret emoticons by focusing on the mouth, which is evident in expressions like :) and :-). In contrast, East Asian users tend to emphasize the eyes, as seen in kaomoji such as (\^ \_ \^). Although the specific study by Yuki et al. (2006) is not fully detailed online, the general concept aligns with findings from cultural psychology: Westerners typically process facial expressions by considering the entire face, particularly the mouth, while East Asians often focus on the eyes.

To further illustrate these cultural differences in emoji perception, the following figure presents a comparative analysis of emoji usage across different cultures, highlighting how context alters their meanings. Note that this figure was created by us and is not attributed to Yuki.


\begin{figure}[h!]
	\centering
	\includegraphics[width=0.8\textwidth]{OurMiller.png}
	\caption{An example of comparative analysis of emoji usage across different cultures (Yuki et al., 2006)}
\end{figure}

\section{The Evolution of Emoji Usage}

Using emoticons in some respects could be considered a language of its own, with each emoticon having a meaning. Like other languages, emoji usage evolved over time and still does to this day. Some of the more recent examples were mentioned in a paper by Levi et al. (\citeyear{Levi2024}), which noted that the face mask emoji was used more frequently during the COVID-19 pandemic and that the meaning of the raised fist emoji is changed to support the Black Lives Matter movement in 2020. Another paper by Robertson et al. (\citeyear{Robertson2021SemanticJQ}) discussed the evolution of emoji usage and the reasons why their meanings changed. These changes were typically caused by trends, specific events, and meme culture. In addition to this, the study mentioned that the usage of emoji could change throughout different seasons. Another study by Barbieri et al.(\citeyear{barbieri2018exploringemojiusageprediction}) noted that emojis like the Christmas tree or snowflake were used more during winter, while emojis related to the sun and camping became more popular during summer.

\section{Sentiment analysis in the Philippines}

Sentiment analysis is an important method of Natural Language Processing (NLP) that measures text sentiment, providing information about people's opinions about diverse subjects, including business, events, trends, and news. Scholars are challenged in the Philippines because Filipino texts have a multilingual background and contain diverse Filipino dialects. Some studies in this context were about sentiment analysis of important issues, such as the 2022 Philippine election (Macrohon et al., 2022) and the Covid-19 vaccine (Villavicencio et al., 2021).

Both used the Naive Bayes technique, a statistical model for predicting sentiment, based on data derived from Twitter. Analyzing tweets about the Philippine elections showed that about 83.90\% of sentiments were negative, 13.49\% were positive, and 2.60\% were neutral (Macrohon et al., 2022). This shows a dominance of negative sentiment about the elections. In comparison, the sentiment analysis about the COVID-19 vaccine showed a better result, with about 83\% of Filipino tweets having positive sentiment (Villavicencio et al., 2021).

Additionally, Campit (2023) did a customer satisfaction study through sentiment analysis of the university president's satisfaction survey. Following previous works, this work used manual annotation and the Naive Bayes algorithm, using Google Translate to overcome the challenge of texts in multiple languages. The study used preprocessing techniques like removing Filipino stop words and stemming. Through manual annotation, 22,380 and 27,834 sentences were positively and neutrally labeled, respectively, with 5,786 negatively labeled, showing a generally positive sentiment about the university president (Campit, 2023). These papers contribute substantially to knowledge of sentiment dynamics in the Filipino case and demonstrate the challenge of analyzing with multilingual and dialectally variant texts.


\begin{center}
	\begin{tabular}{|m{10em} | m{7em} | m{5em} | m{7em} |}
		\hline
		Research Key & Features & Annotation type & Limitations \\ [0.5ex]
		\hline
		FilCon: Filipino Sentiment Lexicon Generation Using Word Level-Annotated Dictionary-Based and (Keen et al., \citeyear{Keen2015}) &
		- Creates a sentiment lexicon
		\newline - Used word level annotation to generate the lexicon
		& Automatic &
		- Was translated to english
		\newline - Focus solely on lexicon construction
		\newline - Low classification accuracy \\
		\hline
		Sentiment Analysis of Tweets on Coronavirus Disease 2019 (COVID-19) Pandemic from Metro Manila, Philippines (Garcia, \citeyear{Garcia2020}) &
		- Lexicon based  based sentiment analysis
		\newline - Used an english lexicon for sentiment analysis
		\newline - Analyzes Filipino sentiment on Covid 19 vaccine &
		Pre labelled &
		- Was translated to english
		\newline - Limited Geographic Scope
		\newline - Data was limited by gathering time frame \\
		\hline
		Application of Naïve Bayes Algorithm in Sentiment Analysis of Filipino, English and Taglish Facebook Comments (Bilog, \citeyear{Bilog2020}) &
		- Machine learning based sentiment analysis
		\newline - Uses Naive Bayes to classify the sentiment of tweets
		\newline - Analyzes facebook comments &
		Manual &
		- Limited sample size
		\newline - Limited diversity in training data
		\newline - Limitations in sentiment classification method \\
		\hline
		Twitter Sentiment Analysis towards COVID-19 Vaccines in the Philippines Using Naïve Bayes (Villavicencio et al., \citeyear{Villavicencio2021}). &
		- Machine learning based sentiment analysis
		\newline - Uses Naive Bayes to classify the sentiment of tweets
		\newline - Analyzes Filipino sentiment on Covid 19 vaccine &
		Manual &
		- Data gathered over a limited time frame
		\newline - Limitations of the free version of RapidMiner application \\
		\hline
	\end{tabular}
\end{center}

\newpage

\begin{center}
	\begin{tabular}{|m{10em} | m{7em} | m{5em} | m{7em} |}
		\hline
		A Semi-Supervised Approach to Sentiment Analysis of Tweets during the 2022 Philippine Presidential Election (Macrohon et al., \citeyear{Macrohon2022}) &
		- Machine learning based sentiment analysis
		\newline - Semi supervised Labeling of tweets
		\newline - Uses Naive Bayes to classify the sentiment of tweets
		\newline - Analyzes the sentiment of the 2022 election &
		Semi Automatic &
		- Data gathered over a limited time frame
		\newline - Dataset imbalance
		\newline - Lack of external validation
		\newline Limited code switching support \\
		\hline
		Analyzing Customer Satisfaction using Support Vector Machine and Naive Bayes Utilizing Filipino Text (Campit, \citeyear{Campit2023}) &
		- Machine learning based sentiment analysis
		\newline -Uses Naive Bayes to classify the sentiment of tweets
		\newline -Analyzes customer sentiments on customer satisfaction &
		Manual &
		- Narrow scope
		\newline - Limited code switching support\\
		\hline
		Affect regression of tweets from the Philippines using emoji-based labels (Co et al., \citeyear{Co2023}) &
		- Machine learning based sentiment analysis
		\newline - Compares CNNs and BERT based models
		\newline - Scores words on based on their emotion &
		Automatic &
		- Bert model was not fine tuned\\
		\hline
	\end{tabular}
\end{center}

\section{Emoji Sentiment Lexicons}

In sentiment analysis there are 2 approaches to classifying sentiment: a machine learning based approach and a rule/lexicon based approach. Sentiment lexicons are a specific set of words, phrases and in our case emojis that carry sentiment that the rule/lexicon based models rely on to assign sentiment scores/labels to text. Sentiment lexicon has 1 or more sentiment scores relating to each word, phrase or emoji. According to a study by Khan et al. (2025) including emojis in sentiment analysis can improve the accuracy of sentiment analysis showing us that emoji sentiment lexicons are relevant in the field of sentiment analysis. Multiple studies have been incorporating emojis lexicons into sentiment analysis to measure sentiment; one such study done by Co et al. (2023) features an emoji sentiment lexicon that they used to help train machine learning models to classify sentiments of tweets.

One of the first emoji sentiment lexicons is included in VADER, a rule based sentiment analysis model with a sentiment lexicon that includes emojis (Hutto \& Gilbert, 2014). Vader has some shortcomings with its representation of emojis as their sentiment scores are based on the words that make up said emoji. Vader was constructed by compiling already existing sentiment lexicons and through crowdsourcing sentiment labels. The annotators labeled each entry with a score from -4 to 4 based on valence (intensity) and were all fluent or native English speakers (Hutto \& Gilbert, 2014).

An essential sentiment lexicon for this study is The Emoji Sentiment Ranking (ESR). This first emoji sentiment lexicon was explicitly made to provide a sentiment map of some of the most frequently used emojis (Novak et al., 2015). The ESR gives each emoji its sentiment scores, which are the sentiment score, its proportion of positive, negative, and neutral labels, and the number of times the emoji appeared in texts. The ESR assigns sentiment scores to each emoji and gives them a sentiment score from 1 (Positive) to -1 (Negative) based on how often the emojis appear in texts with negative and positive labels. The labels used to compile the lexicon were gathered through the manual annotation of many tweets (posts from X, formerly Twitter) made up of 13 different European languages, Novak et al. (2015). This study will use many of the methods mentioned in this study, such as their annotation lexicon construction methods.


\begin{figure}[h!]
	\centering
	\includegraphics[width=0.8\textwidth]{EmojiSentimentLexiconExample.png}
	\caption{An example of an Emoji Sentiment Lexicon}
\end{figure}

\newpage

A lexicon constructed specifically for Arabic was constructed by Alfreihat et al. (2024) which details the creation of an emoji sentiment lexicon based on the Arabic language. Unlike the previous two studies the annotation used for this study was done automatically although a flaw with this study is it translates Arabic text to English before it annotates the texts which potentially fails to capture the certain nuances of the Arabic language. The method this lexicon uses to calculate sentiment scores is similar to the ESR; it makes use of an equation that outputs the average between the positive and negative occurrences of an emoji which results in a sentiment score ranging from -1 to 1 (Alfreihat et al., 2024). The study also details how a sentiment analysis model based on this lexicon can compute the sentiment scores of emojis by getting the sum of the sentiment scores of each emoji that appeared in the piece of text you are classifying.

The previous lexicons mentioned lack dimensions such as feelings and emotions. Sentiment lexicons such as MLE besides just scoring valence (positivity and negativity) also give each emoji multiple emotion scores (Godard \& Holtzman, 2022) which allows deeper understanding by providing an additional dimension that being emotion besides just valence. It does this by rating emojis based on valence and 8 other human emotions such as anger, anticipation, disgust, fear, joy, sadness, surprise, and trust.


\begin{figure}[h!]
	\centering
	\includegraphics[width=0.8\textwidth]{GodardAndHoltzmanSL.png}
	\caption{Godard \& Holtzman's (2022) Emoji Sentiment Lexicon}
\end{figure}

\newpage

This image shows the top emojis in each emotion category in the Multidimensional emoji of lexicons (Godard \& Holtzman, 2022). This lexicon was created through leveraging another multidimensional lexicon, specifically the NRC emotion lexicon, to assign sentiment labels to tweets containing emojis. Another sentiment lexicon known as Lisbon Emoji and Emoticon Database (LEED) scores emojis based on feelings such as aesthetic appeal, familiarity, visual complexity, concreteness, valence, arousal, and meaningfulness. (Rodrigues et al., 2017) Unlike the previous lexicon instead of annotating a large number of texts with emojis this lexicon asked its annotators to rate the emojis directly and averaging out the scores they were given.


\begin{center}
	\begin{tabular}{|m{10em} | m{7em} | m{5em} | m{7em} |}
		\hline
		Research & Key features & Annotation type & Limitations \\
		\hline
		Sentiment of Emojis (Novak et al, \citeyear{KraljNovak2015}) &
		- Rates emojis based on valence
		\newline - Supports 13 european languages &
		Manual &
		- Only European languages were used
		\newline - Emotion was not fully captured with \\
		\hline
		Lisbon Emoji and Emoticon Database (LEED): Norms for emoji and emoticons in seven evaluative dimensions (Rodrigues et al., \citeyear{Rodrigues2017}) &
		- Rates emojis based on aesthetic appeal, familiarity, visual complexity, concreteness, valence, arousal, and meaningfulness
		\newline - Supports Portuguese 
		\newline - Was created by annotating emojis directly &
		Manual &
		- Data was only gathered from Portuguese participants \\
		\hline
		The Multidimensional Lexicon of Emojis: A New Tool to Assess the Emotional Content of Emojis (Godard \& Holtzman, \citeyear{Godard2022}) &
		- Rates emojis based on valence, anger, anticipation, disgust, fear, joy, sadness, surprise, and trust
		\newline - Supports English 
		\newline - Was created by annotating tweets containing emojis using the NRC lexicon &
		Automatic & 
		- Emotional nuances such as sarcasm was not accounted for
		\newline - Temporal and cultural factors were not considered
		Mainly English content \\
		\hline
		Emo-SL Framework: Emoji Sentiment Lexicon Using Text-Based Features and Machine Learning for Sentiment Analysis Manar Alfreihat, et al. (\citeyear{Alfreihat2024}) &
		- Rates emojis based on valence
		\newline - Supports Arabic &
		Automatic &
		- Translated Arabic texts to English not capturing their full meaning. Colloquial expressions and dialects not fully captured \\
		\hline
	\end{tabular}
	\newpage
	\begin{tabular}{|m{10em} | m{7em} | m{5em} | m{7em} |}
		\hline
		Emojis in Lexicon-Based Sentiment Analysis: Creating Emoji Sentiment Lexicons from Unlabeled Corpora (Haak, \citeyear{Haak2021}) &
		- Supports non UTF-8 emojis from twitch &
		Automatic &
		- Emotion was not fully captured with
		\newline - Temporal factors were not considered \\
		\hline
		Automatic Construction of an Emoji Sentiment Lexicon (Mayu Kimura \& Marie Katsurai, \citeyear{Kimura_Katsurai2017}) &
		- Rates emojis based on valence, anger, sadness, fear, disgust, and happiness &
		Automatic &
		- Only has 236 emojis 
		\newline - Does not capture other positive emotions such as anticipation surprise and trust \\
		\hline
	\end{tabular}
\end{center}

As emojis have a non-insignificant effect on the outcome of sentiment analysis tasks, fully understanding how emojis may affect such outcomes is important. One research suggests, "Emoji are an important substitute for non-verbal cues (such as facial expressions) in online written communication" (Chen et al., 2024). Online, emojis' meaning varies by age, gender, and especially culture (Chen et. al., 2024). However, in practice, a given emoji may not carry the same emotional nuance for a Filipino user as it does for Western users. For example, a study by Togans et. al. in 2021 regarding cross-cultural communication found that East Asian Participants (e.g., Chinese and Korean) use emojis more frequently and carefully than Americans, reflecting indirect communication styles and concern for social harmony. These cultural patterns have practical effects. For instance, Emojis that indicate a smile or laughter in one culture may carry additional connotations in another culture.

Most existing sentiment analysis tools and lexicons for emojis assume a Western baseline. For example, Novak et al.'s Emoji Sentiment Ranking (ESR) was derived from over 70,000 tweets in 13 European languages, producing sentiment scores for 751 emojis based on primarily Western usage. Similarly, general-purpose English lexicons (e.g., NRC EmoLex) and rule-based analyzers like VADER are tuned to Western-style emoticons and emoji use. When applied to Filipino texts, these tools can produce systematic errors. Ljubešić and Fišer, in a 2016 study, note that simpler models and Western lexicons may not capture local usage patterns, recommending instead to merge emoji features with more contextual embeddings for better results.

As a result of this, cultural biases have been directly documented. In an Arabic context, Hakami et. al. in 2021 found that applying the European emoji lexicon produced inaccuracies; while basic emojis aligned in both Arabic and European lexicons, many others showed divergent sentiment ratings. Tools like VADER or generic lexicons do not effectively capture these differences. For example, VADER's rule-based scoring does not distinguish between an emoji's use in a Western versus Filipino context and hence can misclassify the sentiment of emoji-rich Filipino messages. Empirical results observed by Alfreihat et. al. 2024 support the need for adaptation: an Arabic emoji sentiment lexicon integrated into a classifier raised accuracy dramatically (F1=89\%) and outperformed VADER by a large margin. In that work, adding an Arabic-specific emoji lexicon to machine learning features improved accuracy by ~26.7\% over text-only baselines. Broadly, Lexicons underscore that culturally adapted emoji lexicons substantially improve sentiment analysis performance.



\section{Language Models}

Language models were powerful tools that automated text classification and filtering. Language models could be used to help create sentiment lexicons. They could classify emotions and sentiments in texts, which helped in the process of sentiment analysis. An example of a language model was Word2Vec, an algorithm used for Natural Language Processing (NLP) that employed neural networks to learn word embeddings from large texts (\citeyear{Jang2019}). It understood and vectorized words, based on the hypothesis that words with similar meanings should be close in distance. Word2Vec had multiple models, such as CBOW, Skip-gram, GloVe, and FastText. The study by Jang et al. (\citeyear{Jang2019}) only discussed CBOW, which predicted words based on context, and Skip-gram, which predicted context based on words. Another study by Jiang et al. (\citeyear{Jiang2021}) sought to alleviate the lack of labeled resources for Tagalog by developing three pre-trained language models (PLMs). BERT, ELECTRA, and RoBERTa. BERT, which stands for Bidirectional Encoder Representations for Transformers, learned from bidirectional representations from unlabeled texts and gained context by analyzing texts both forward and backward. BERT consisted of two unsupervised tasks: Next-Sentence Prediction (NSP) and Masked Language Modeling (MLM). NSP determined whether two sentences were logically connected, while MLM masked some words in the input and predicted what those masked words were. RoBERTa, another PLM, was based on BERT but was trained with more data, removed NSP, and added dynamic word masking. ELECTRA, the final PLM, used a discriminator and a generator and employed a pre-training task called replaced token detection.




\begin{comment}
	%
	% IPR acknowledgement: the contents within this comment are from Ethel Ong's slides on RRL.
	%
	Guide on Writing your Related Works chapter
	
	1. Identify the keywords with respect to your research
	One keyword = One document section
	Examples: 2.1 Story Generation Systems
	2.2 Knowledge Representation
	
	2.  Find references using these keywords
	
	3.  For each of the references that you find,
	Check: Is it relevant to your research?
	Use their references to find more relevant works.
	
	4. Identify a set of criteria for comparison.
	It will serve as a guide to help you focus on what to look for
	
	5. Write a summary focusing on -
	What: A short description of the work
	How: A summary of the approach it utilized
	Findings: If applicable, provide the results
	Why: Relevance to your work
	
	6. At the end of each section,  show a Table of Comparison of the related works 
	and your proposed project/system
	
\end{comment}